\section{面向对象（上）}

\begin{question}
	以下关于类和对象的说法正确的是
	\begin{options}
		\item 类是对象的实例
		\item 对象是类的模板
		\item 一个类只能创建一个对象
		\item 对象是类的具体实例
	\end{options}
	\begin{answer}
		D
	\end{answer}
	\begin{explanation}
		类是抽象模板，对象是类的具体实例。一个类可以创建多个对象。
	\end{explanation}
\end{question}

\begin{question}
	以下哪个是Java中的对象引用传递？
	\begin{options}
		\item 基本数据类型参数传递
		\item 对象作为参数传递
		\item 数组作为参数传递
		\item 以上都是
	\end{options}
	\begin{answer}
		B
	\end{answer}
	\begin{explanation}
		对象作为参数传递时传递的是引用（地址），属于引用传递。
	\end{explanation}
\end{question}
